\section{Prima terkait}

Sekarang kita akan memulai teori struktur modul noetherian. Langkah pertama
adalah mengaitkan suatu kumpulan ideal prima kepada setiap modul, yang disebut
\emph{prima terkait}; kumpulan ini termuat dalam kumpulan ideal prima yang lebih besar (yaitu
\emph{dukungan}), tempat
pelokalan modul tersebut tak nol.

\subsection{Dukungan}
 Misalkan $R$ suatu gelanggang noetherian. Suatu modul-$R$, $M$, hendaknya dipandang
 sebagai sesuatu yang menyerupai berkas vektor, yang dengan suatu cara
tersebar di atas ruang topologis $\spec R$. Jika $\mathfrak{p} \in \spec R$, tetapkan
\( \k(\mathfrak{p}) = \mathrm{medan \ pecahan \ } R/\mathfrak{p}  ,\)
yang merupakan medan residu dari $R_{\mathfrak{p}}$. Jika $M$ sebarang modul-$R$, kita
dapat mempertimbangkan $M \otimes_R \k(\mathfrak{p})$ untuk setiap $\mathfrak{p}$; objek ini merupakan
ruang vektor di atas $\k(\mathfrak{p})$. Jika $M$ terbangkitkan hingga, maka $M \otimes_R
\k(\mathfrak{p})$ merupakan ruang vektor berdimensi hingga.

\begin{definition} 
Misalkan $M$ suatu modul-$R$ terbangkitkan hingga. Maka $\supp M$, yaitu
\textbf{dukungan} dari $M$, didefinisikan sebagai himpunan ideal prima
$\mathfrak{p} \in \spec R$ sedemikian sehingga
\( M \otimes_R \k(\mathfrak{p}) \neq 0.  \)
\end{definition} 

Kita hendaknya memandang suatu modul $M$ sebagai sesuatu yang menyerupai berkas vektor
di atas ruang topologis
$\spec R$. Kepada setiap $\mathfrak{p} \in \spec R$, kita kaitkan ruang vektor $M
\otimes_R \k(\mathfrak{p})$; inilah “seratnya”. Tentu saja, intuisi bahwa
$M$ merupakan berkas vektor agak terbatas, sebab serat-seratnya
pada umumnya tidak mempunyai dimensi yang sama.
Meskipun demikian, kita dapat membicarakan dukungan, yakni himpunan titik tempat
“serat” tersebut tak nol.

Perhatikan bahwa $\mathfrak{p} \in \supp M$ jika dan hanya jika $M_{\mathfrak{p}} \neq 0$. Hal ini
karena
\[ (M \otimes_R R_{\mathfrak{p}})/( \mathfrak{p} R_{\mathfrak{p}} (M \otimes_R
R_{\mathfrak{p}}))  = M_{\mathfrak{p}}
\otimes_{R_{\mathfrak{p}}} \k(\mathfrak{p})  \]
dan kita dapat menggunakan Lema Nakayama di atas gelanggang lokal $R_{\mathfrak{p}}$. (Di sini kita
menggunakan fakta bahwa $M$ terbangkitkan hingga.)

Berkas vektor yang dukungannya kosong adalah nol. Karena itu, proposisi mudah
berikut bersifat intuitif:

\begin{proposition} 
$M = 0$ jika dan hanya jika $\supp M = \emptyset$. 
\end{proposition}
\begin{proof} 
Memang, $M=0$ jika dan hanya jika $M_{\mathfrak{p}} = 0$ untuk semua ideal prima
$\mathfrak{p} \in \spec R$. Hal ini setara dengan $\supp M = \emptyset$.
\end{proof} 
 
\begin{exercise} 
Misalkan $0 \to M' \to M \to M'' \to 0$ eksak. Maka 
\[ \supp M = \supp M' \cup \supp M''.  \]
\end{exercise} 


Kita akan segera melihat bahwa $\supp M$ tertutup di dalam $\spec R$. Dalam suatu
pengertian, kita membayangkan bahwa $M$ hidup pada himpunan bagian tertutup $\supp M$ ini.



\subsection{Prima terkait}
Di seluruh bagian ini, $R$ adalah gelanggang noetherian. \emph{Prima
terkait} dari suatu modul $M$ adalah ideal-ideal prima yang muncul sebagai anihilator
elemen-elemen dalam $M$. Seperti akan kita lihat, dukungan suatu modul ditentukan oleh
prima terkaitnya. Secara khusus, prima terkait memuat “titik-titik generik”
(yakni ideal-ideal prima minimal) dari dukungan tersebut. Namun, dalam beberapa kasus,
prima terkait dapat memuat lebih banyak ideal. 

\begin{quote}\small\textbf{Jembatan edisi (asli).}
Untuk $x\in M$, tuliskan
$\ann_R(x)=\{r\in R:rx=0\}$ untuk anihilator $x$; jika gelanggangnya
sudah jelas, notasi ini disingkat menjadi $\ann(x)$.\end{quote}

\begin{definition} 
Misalkan $M$ suatu modul-$R$ terbangkitkan hingga. Ideal prima $\mathfrak{p}$ dikatakan
\textbf{terkait} dengan $M$ jika terdapat elemen $x \in M$ sedemikian sehingga
$\mathfrak{p}$ merupakan anihilator $x$. Himpunan prima terkait ditulis
$\ass(M)$.
\end{definition} 

Perhatikan bahwa anihilator suatu elemen $x \in M$ belum tentu prima, tetapi
mungkin saja anihilator tersebut prima; dalam hal itu, anihilator tersebut
terkait.

\begin{exercise} 
Tunjukkan bahwa $\mathfrak{p} \in \ass(M)$ jika dan hanya jika terdapat penyisipan
$R/\mathfrak{p} \hookrightarrow M$.
\end{exercise} 

\begin{exercise} 
Misalkan $\mathfrak{p} \in \spec R$. Maka $\ass(R/\mathfrak{p}) =
\left\{\mathfrak{p}\right\}$.
\end{exercise} 

\begin{example}
Ambil $R=k[x,y,z]$, dengan $k$ suatu domain integral, dan misalkan $I = (x^2-yz,x(z-1))$. Setiap
 ideal prima yang terkait dengan $I$ harus memuat $I$, jadi mari kita pertimbangkan
   $\mathfrak{p}=(x^2-yz,z-1)=(x^2-y,z-1)$, yang sama dengan $I:x$. Ideal ini prima karena $R/\mathfrak{p} = k[x]$,
   yang merupakan domain. Untuk melihat bahwa $(I:x)\subset \mathfrak{p}$, andaikan $tx\in I\subset \mathfrak{p}$; karena
   $x\not\in \mathfrak{p}$, diperoleh $t\in p$, seperti yang dikehendaki.

   Terdapat dua prima terkait lainnya, tetapi kita tidak akan mencarinya di sini.
 \end{example}


Kita mulai dengan membuktikan bahwa $\ass(M) \neq \emptyset$ untuk modul tak nol. 
\begin{proposition} \label{assmnonempty} 
Jika $M \neq 0$, maka $M$ mempunyai prima terkait.
\end{proposition} 
\begin{proof}  Pertimbangkan kumpulan ideal dalam $R$ yang muncul sebagai
anihilator suatu elemen tak nol dalam $M$. 
Misalkan $I \subset R$ elemen maksimal dari kumpulan tersebut. Keberadaan $I$ dijamin oleh sifat noetherian
$R$.
Maka $I = \ann(x)$ untuk suatu $x \in M$, sehingga $1 \notin I$ karena anihilator suatu elemen tak nol bukan seluruh
gelanggang.

Kita mengklaim bahwa
$I$ prima, sehingga $I \in \ass(M)$.  
Memang, andaikan $ab \in I$ dengan $a,b \in R$. Ini berarti
\[ (ab)x = 0.  \]
Pertimbangkan anihilator $\ann(bx)$ dari $bx$. Anihilator ini memuat anihilator $x$, yaitu $I$;
anihilator tersebut juga memuat $a$.

Ada dua kasus. Jika $bx = 0$, maka $ b \in I$ dan pembuktian selesai. Andaikan
sebaliknya bahwa $bx \neq 0$. Dalam kasus ini, $\ann(bx)$ memuat $(a) + I$, yang
 memuat $I$. Berdasarkan kemaksimalan, harus berlaku $\ann(bx) = I$ dan $ a \in
 I$. 

 Dalam kedua kasus tersebut, kita mendapati bahwa salah satu dari $a,b $ termasuk dalam $I$, sehingga $I$
 prima. 

\end{proof} 

\begin{example}[Modul tanpa prima terkait]
Tanpa hipotesis noetherian, \rref{assmnonempty}
\emph{salah}. Pertimbangkan $R = \mathbb{C}[x_1, x_2, \dots]$, gelanggang polinomial
di atas $\mathbb{C}$ dalam tak hingga banyak variabel, dan ideal $I = (x_1,
x_2^2, x_3^3, \dots) \subset R$.
Klaimnya adalah 
\[ \ass(R/I ) = \emptyset.  \]
Untuk melihatnya, andaikan suatu ideal prima $\mathfrak{p}$ merupakan anihilator dari
$\overline{f}\in R/I$. Maka $\overline{f}$ dapat diangkat menjadi $f \in R$; akibatnya,
$\mathfrak{p}$ tepat merupakan himpunan semua $g \in R$ sedemikian sehingga $fg \in I$.
Sekarang $f$ hanya memuat sejumlah hingga variabel $x_i$, katakanlah $x_1, \dots,
x_n$. Jelas bahwa $x_{n+1}^{n+1} f \in I$ (sehingga $x_{n+1}^{n+1} \in
\mathfrak{p}$), tetapi $x_{n+1} f \notin I$ (sehingga $x_{n+1} \notin \mathfrak{p}$).
Jadi, $\mathfrak{p}$ bukan ideal prima, suatu kontradiksi.
\end{example} 

Sekarang kita akan menunjukkan bahwa banyaknya prima terkait berhingga.

\begin{proposition} \label{finiteassm}
Jika $M$ terbangkitkan hingga, maka $\ass(M)$ berhingga.
\end{proposition}

Gagasannya ialah menggunakan fakta bahwa $M$ terbangkitkan hingga untuk menyusun
$M$ dari sejumlah hingga bagian, lalu menggunakannya untuk membatasi banyaknya prima
terkait pada setiap bagian. Untuk itu kita memerlukan:

\begin{lemma} \label{assexact} 
Andaikan terdapat barisan eksak modul-$R$ terbangkitkan hingga
\[ 0 \to M' \to M \to M'' \to 0.  \]
Maka 
\[\ass(M') \subset \ass(M) \subset \ass(M') \cup \ass(M'')  \]
\end{lemma} 
\begin{proof} 
Klaim pertama sudah jelas. Jika $\mathfrak{p}$ merupakan anihilator
suatu $x \in M'$, maka ia juga merupakan anihilator sesuatu dalam $M$ (yaitu citra
$x$), karena
$M' \to M$ injektif. Jadi, $\ass(M') \subset \ass(M)$. 

Arah yang lebih sulit adalah penyertaan lainnya. Andaikan $\mathfrak{p} \in \ass(M)$.
Maka terdapat $x \in M$ sedemikian sehingga
$\mathfrak{p} = \ann(x).$
Pertimbangkan submodul $Rx \subset M$. Jika $Rx \cap M' \neq 0$, kita dapat
memilih $y \in Rx \cap M' - \left\{0\right\}$. Kita mengklaim bahwa $\ann(y) =
\mathfrak{p}$, sehingga $\mathfrak{p} \in \ass(M')$.
Untuk melihatnya, $ y = ax$ untuk suatu $a \in R$. Anihilator $y$ adalah himpunan elemen
$b \in R$ sedemikian sehingga
\[ abx = 0  \]
atau, secara ekuivalen, himpunan semua $b \in R$ sedemikian sehingga $ab \in \mathfrak{p} =
\ann(x)$. Namun, $y = ax \neq 0$, sehingga $a \notin \mathfrak{p}$. Oleh
karena itu, syarat $b \in \ann(y)$ sama dengan $b \in \mathfrak{p}$. Dengan
kata lain, 
\[ \ann(y) = \mathfrak{p}  \]
yang membuktikan klaim tersebut.

Sekarang andaikan $Rx \cap M' = 0$. Misalkan $\phi: M \twoheadrightarrow M''$
pemetaan surjektif. Kita mengklaim bahwa $\mathfrak{p} = \ann(\phi(x))$, sehingga
akibatnya
$\mathfrak{p} \in \ass(M'')$. Pembuktiannya sebagai berikut. Jelas bahwa $\mathfrak{p}$
menganihilasi $\phi(x)$ karena ia menganihilasi $x$. Andaikan $a \in \ann(\phi(x))$.
Ini berarti $\phi(ax) = 0$, sehingga $ax \in \ker \phi=M'$; tetapi $\ker \phi \cap Rx =
0$. Jadi, $ax = 0$, dan akibatnya $a \in \mathfrak{p}$. Maka $\ann(\phi(x)) = \mathfrak{p}$. 
\end{proof} 

Langkah berikutnya dalam pembuktian \rref{finiteassm} adalah bahwa setiap
modul terbangkitkan hingga
memiliki suatu filtrasi yang setiap hasil bagi berturutannya berbentuk sangat baik.
Hasil ini cukup berguna dan akan dirujuk lagi kelak.

\begin{proposition}[D{\'e}vissage] \label{filtrationlemma} \label{devissage}
Untuk setiap modul-$R$ terbangkitkan hingga $M$, terdapat suatu filtrasi hingga
\[ 0 = M_0 \subset M_1 \subset \dots \subset M_k = M  \]
sedemikian sehingga hasil-hasil bagi berturutan $M_{i+1}/M_i$ isomorfik dengan berbagai
$R/\mathfrak{p}_i$, dengan $\mathfrak{p}_i \subset R$ ideal prima.
\end{proposition} 
\begin{proof} 
Misalkan $M' \subset M$ maksimal di antara submodul-submodul yang memiliki filtrasi semacam itu
(yang berakhir pada $M'$).
Kita ingin menunjukkan bahwa $M' = M$. Submodul $M'$ terdefinisi dengan baik karena
$0$ mempunyai filtrasi semacam itu dan $M$
noetherian.  

Terdapat suatu filtrasi
\[ 0 = M_0 \subset M_1 \subset \dots \subset M_l = M' \subset M  \]
dengan hasil-hasil bagi berturutan, \emph{kecuali} mungkin yang terakhir $M/M'$, berbentuk
$R/\mathfrak{p}_i $ untuk $\mathfrak{p}_i \in \spec R$.
Jika $M' = M$, pembuktian selesai. Jika tidak, pertimbangkan
hasil bagi $M/M' \neq 0$. Modul $M/M'$ mempunyai prima terkait. Jadi, terdapat
ideal prima $\mathfrak{p}$ yang merupakan anihilator suatu $x \in M/M'$. Ini berarti
terdapat penyisipan 
\[ R/\mathfrak{p} \hookrightarrow M/M'.  \]
Sekarang, ambil $M_{l+1}$ sebagai prapeta di dalam $M$
dari $R/\mathfrak{p} \subset M/M'$. 
Kemudian kita dapat mempertimbangkan filtrasi hingga
\[ 0 = M_0 \subset M_1 \subset \dots \subset M_{l+1} , \]
yang seluruh hasil bagi berturutannya berbentuk $R/\mathfrak{p}_i$; hal ini
karena $M_{l+1}/M_l = M_{l+1}/M'$ mempunyai bentuk tersebut berdasarkan konstruksi.
Dengan demikian, kita telah memperpanjang filtrasi ini satu
langkah lagi, suatu kontradiksi karena
$M'$ diasumsikan maksimal.
\end{proof} 

Sekarang kita siap mencapai tujuan dan membuktikan bahwa $\ass(M)$
selalu merupakan himpunan berhingga. 
\begin{proof}[Pembuktian \rref{finiteassm}]
Andaikan $M$ terbangkitkan hingga. Ambil filtrasi kita
\[ 0 = M_0 \subset M_1 \subset \dots \subset M_k = M.  \]
Dengan induksi, kita tunjukkan bahwa $\ass(M_i)$ berhingga untuk setiap $i$. Hal ini jelas
benar untuk $i=0$. Sekarang andaikan $\ass(M_i)$ berhingga; kita membuktikan hal yang sama untuk
$\ass(M_{i+1})$. Kita mempunyai barisan eksak
\[ 0 \to M_i \to M_{i+1} \to R/\mathfrak{p}_i \to 0  \]
yang, berdasarkan \rref{assexact}, menyiratkan
\[ \ass(M_{i+1}) \subset \ass(M_i) \cup \ass(R/\mathfrak{p}_i) = \ass(M_i)
\cup \left\{\mathfrak{p}_i\right\} , \]
sehingga $\ass(M_{i+1})$ juga berhingga.
Dengan induksi, kini jelas bahwa $\ass(M_i)$ berhingga untuk setiap $i$.

Ini membuktikan proposisi tersebut; pembuktian ini juga menunjukkan bahwa banyaknya
prima terkait tidak melebihi panjang filtrasi. 
\end{proof} 


Terakhir, kita mencirikan pembagi nol pada $M$ melalui prima terkait.
Pencirian terakhir dari hasil ini akan berguna kelak.
Sebagai contoh, hasil tersebut menyiratkan bahwa jika $R$ lokal dan $\mathfrak{m}$ ideal
maksimalnya, maka jika setiap elemen $\mathfrak{m}$ merupakan pembagi nol pada modul
terbangkitkan hingga
$M$, berlaku $\mathfrak{m} \in \ass(M)$.

\begin{proposition} \label{assmdichotomy}
Jika $M$ suatu modul terbangkitkan hingga di atas gelanggang noetherian $R$, maka semua
pembagi nol pada $M$ adalah gabungan $\bigcup_{\mathfrak{p} \in \ass(M)}
\mathfrak{p}$.

Lebih kuat lagi, jika $I \subset R$ sebarang ideal yang terdiri atas pembagi nol pada
$M$, maka $I$ termuat dalam suatu prima terkait.
\end{proposition} 
\begin{proof} 
Setiap prima terkait merupakan anihilator suatu elemen $M$, sehingga terdiri
atas pembagi nol. Sebaliknya, jika $a \in R$ menganihilasi $x \in M$, maka $a$
termasuk dalam setiap prima terkait dari modul tak nol $Ra \subset M$. (Setidaknya
ada satu berdasarkan \cref{finiteassm}.)

Untuk pernyataan terakhir, kita menggunakan penghindaran prima (\cref{primeavoidance}): jika $I$ terdiri atas
pembagi nol, maka berdasarkan bagian pertama pembuktian, $I$ termuat dalam gabungan $\bigcup_{\mathfrak{p} \in \ass(M)}
\mathfrak{p}$. Ini merupakan gabungan berhingga berdasarkan
\cref{assmfinite}, sehingga penghindaran prima menyiratkan bahwa $I$ termuat dalam salah satu
ideal prima tersebut.
\end{proof} 


\begin{exercise} 
Untuk setiap modul $M$ di atas sebarang gelanggang $R$ (yang tidak harus noetherian),
himpunan pembagi nol-$M$, $\mathcal{Z}(M)$, merupakan gabungan ideal-ideal prima. Secara umum, terdapat
 pencirian mudah bagi himpunan $Z$ yang merupakan gabungan ideal-ideal prima: tepat ketika
 $R\smallsetminus Z$ merupakan \emph{himpunan multiplikatif jenuh}. Ini adalah
 Teorema 2 Kaplansky.
 \begin{definition}
   Suatu himpunan multiplikatif $S\neq \varnothing$ disebut \emph{himpunan multiplikatif jenuh} jika
   untuk semua $a,b\in R$, berlaku $a,b\in S$ jika dan hanya jika $ab\in S$. (“Himpunan multiplikatif” hanya
   berarti bagian “jika”.)
 \end{definition}
 Untuk melihat bahwa $\mathcal{Z}(M)$ merupakan gabungan ideal-ideal prima, cukup periksa bahwa komplemennya adalah
 himpunan multiplikatif jenuh.
\end{exercise}

\subsection{Pelokalan dan $\ass(M)$}

Ternyata sangat memudahkan bahwa konstruksi $M  \to \ass(M)$
berperilaku sebaik yang dapat kita harapkan terhadap
pelokalan. Bahkan, hal ini memungkinkan kita mereduksi argumen ke kasus gelanggang lokal,
yang jauh lebih sederhana.

Seperti biasa, misalkan $R $ noetherian dan $M$ suatu modul-$R$ terbangkitkan hingga. 
Misalkan pula $S \subset R$ suatu himpunan bagian multiplikatif. 
Maka $S^{-1}M$ merupakan modul terbangkitkan hingga di atas gelanggang noetherian
$S^{-1}M$. Jadi, masuk akal untuk mempertimbangkan $\ass(M) \subset \spec R$ dan
$\ass(S^{-1}M) \subset \spec S^{-1}R$. Namun, kita juga mengetahui bahwa $\spec S^{-1}R
\subset \spec R$ hanyalah himpunan ideal prima $R$ yang tidak beririsan dengan $S$.
Dengan demikian, kita dapat membandingkan langsung $\ass(M)$ dan $\ass(S^{-1}M)$, serta dapat
menduga (dan ternyata benar) bahwa $\ass(S^{-1}M) = \ass(M) \cap \spec
S^{-1}R$.
\begin{proposition} \label{assmlocalization}
Misalkan $R$ noetherian, $M$ terbangkitkan hingga, dan $S \subset R$ tertutup secara multiplikatif. 
Maka 
\[ \ass(S^{-1}M)  = \left\{S^{-1}\mathfrak{p}: \mathfrak{p} \in \ass(M),
\mathfrak{p}\cap S  = \emptyset \right\} . \]
\end{proposition} 
\begin{proof} 
Pertama-tama kita membuktikan arah yang mudah, yakni bahwa $\ass(S^{-1}M)$
\emph{memuat} ideal-ideal prima dalam $\spec S^{-1}R \cap \ass(M)$. 

Andaikan $\mathfrak{p} \in \ass(M)$ dan
$\mathfrak{p} \cap S = \emptyset$. Maka $\mathfrak{p} = \ann(x)$ untuk suatu $x
\in M$. Anihilator $x/1 \in S^{-1}M$ tepat sama dengan $S^{-1}\mathfrak{p}$, seperti dapat
diperiksa langsung. Jadi, $S^{-1}\mathfrak{p} \in \ass(S^{-1}M)$.
Dengan demikian, kita memperoleh penyertaan yang mudah.

Sekarang kita buktikan penyertaan yang lebih sulit.  
Tuliskan pemetaan pelokalan $R \to S^{-1}R  $ sebagai $\phi$.
Misalkan $\mathfrak{q} \in \ass(S^{-1}M)$. Berdasarkan definisi, ini berarti $\mathfrak{q} =
\ann(x/s)$ untuk suatu $x \in M$, $s \in S$. Kita ingin melihat bahwa
$\phi^{-1}(\mathfrak{q}) \in \ass(M) \subset \spec R$.
Berdasarkan definisi, $\phi^{-1}(\mathfrak{q})$ adalah himpunan elemen $a \in R$ sedemikian sehingga
\[ \frac{ax}{s} = 0 \in S^{-1}M . \]
Dengan kata lain, berdasarkan definisi pelokalan, himpunan ini adalah 
\[  \phi^{-1}(\mathfrak{q}) = \bigcup_{t \in S} \left\{a \in R: atx = 0 \in M\right\} = \bigcup \ann(tx)
\subset R.\]
Namun, kita mengetahui bahwa di antara elemen-elemen berbentuk $\ann(tx)$, terdapat
elemen \emph{maksimal} $I=\ann(t_0 x)$ untuk suatu $t_0 \in S$, sebab $R$
noetherian. Klaimnya adalah $I = \phi^{-1}(\mathfrak{q})$, sehingga
$\phi^{-1}(\mathfrak{q}) \in \ass(M)$.

Memang, setiap anihilator lain $I' = \ann(tx)$ (untuk $t \in S$) harus termuat dalam $\ann(t_0 t x)$. Akan tetapi, 
\( I \subset \ann(t_0 x)  \)
dan $I$ maksimal, sehingga $I = \ann(t_0 t x)$ dan
\( I' \subset I.  \) Dengan kata lain, $I$ memuat semua anihilator lainnya
$\ann(tx)$ untuk $t \in S$.
Secara khusus, gabungan besar di atas, yakni $\phi^{-1}(\mathfrak{q})$, tepat sama dengan
\( I = \ann(t_0 x).  \)
Jadi, $\phi^{-1}(\mathfrak{q}) = \ann(t_0x)$ termasuk dalam $\ass(M)$.
Ini berarti bahwa setiap prima terkait
dari $S^{-1}M$ berasal dari prima terkait $M$, dan pembuktian selesai.
\end{proof} 



\begin{exercise} 
Tunjukkan bahwa jika $M$ suatu modul terbangkitkan hingga di atas gelanggang noetherian, maka
pemetaan
\[ M \to \bigoplus_{\mathfrak{p} \in \ass(M)} M_{\mathfrak{p}}  \]
injektif. Apakah hal ini benar jika $M$ tidak terbangkitkan hingga?
\end{exercise} 

\subsection{Prima terkait menentukan dukungan}
Klaim berikutnya adalah bahwa dukungan dan prima terkait saling berhubungan.

\begin{proposition}\label{supportassociated} Dukungan adalah penutupan dari prima terkait:
\[ \supp M  = \bigcup_{\mathfrak{q} \in \ass(M)}
\overline{\left\{\mathfrak{q}\right\}} \]
\end{proposition} 

Berdasarkan definisi topologi Zariski, ini berarti bahwa suatu ideal prima $\mathfrak{p}
\in \spec R$ termasuk dalam $\supp M$ jika dan hanya jika ideal itu memuat prima
terkait. 

\begin{proof} 
Pertama, kita menunjukkan bahwa $\supp(M)$ memuat himpunan ideal prima
$\mathfrak{p} \in \spec R$ yang memuat suatu prima terkait; ini akan menyiratkan
bahwa $\supp(M) \supset \bigcup_{\mathfrak{q} \in \ass(M)}
\overline{\left\{\mathfrak{q}\right\}}$. Jadi, misalkan $\mathfrak{q}$ suatu
prima terkait dan $\mathfrak{p} \supset \mathfrak{q}$. Kita harus menunjukkan bahwa
\[ \mathfrak{p} \in \supp M, \ \text{yakni} \ M_{\mathfrak{p}} \neq 0.  \]
Namun, karena $\mathfrak{q} \in \ass(M)$, terdapat pemetaan injektif
\[ R/\mathfrak{q} \hookrightarrow M , \]
sehingga pelokalan memberikan pemetaan injektif
\[ (R/\mathfrak{q})_{\mathfrak{p}} \hookrightarrow M_{\mathfrak{p}}.  \]
Di sini, objek pertama $(R/\mathfrak{q})_{\mathfrak{p}}$ tak nol karena tidak ada elemen tak nol dalam $R/\mathfrak{q}$ yang dapat
dianihilasi oleh sesuatu di luar $\mathfrak{p}$. Jadi, $M_{\mathfrak{p}} \neq
0$, dan $\mathfrak{p} \in \supp M$. 

Sekarang kita buktikan penyertaan sebaliknya. Andaikan $\mathfrak{p} \in \supp M$. Kita
harus menunjukkan bahwa $\mathfrak{p}$ memuat suatu prima terkait.  
Berdasarkan asumsi, $M_{\mathfrak{p}} \neq 0$, dan $M_{\mathfrak{p}}$ merupakan modul
terbangkitkan hingga di atas gelanggang noetherian $R_{\mathfrak{p}}$. Jadi, $M_{\mathfrak{p}}$ mempunyai
prima terkait.
Berdasarkan \rref{assmlocalization}, diperoleh bahwa $\ass(M) \cap \spec
R_{\mathfrak{p}}$ tak kosong. Karena ideal-ideal prima $R_{\mathfrak{p}}$
berkorespondensi dengan ideal-ideal prima yang termuat dalam $\mathfrak{p}$, terdapat
ideal prima yang termuat dalam $\mathfrak{p}$ dan termasuk dalam $\ass(M)$. Inilah
tepatnya yang ingin kita buktikan.
\end{proof} 


\begin{corollary} \label{suppisclosed} Untuk $M$ terbangkitkan hingga,  
$\supp M$ tertutup. Lebih lanjut, setiap elemen minimal dari $\supp M$ termasuk dalam
$\ass(M)$.
\end{corollary} 
\begin{proof} 
Memang, hasil di atas menyatakan bahwa
\[ \supp M  = \bigcup_{\mathfrak{q} \in \ass(M)}
\overline{\left\{\mathfrak{q}\right\}}. \]
Karena $\ass(M)$ berhingga, diperoleh bahwa $\supp M$ tertutup.
Persamaan di atas juga menunjukkan bahwa setiap elemen minimal dari $\supp M$ harus merupakan
prima terkait.
\end{proof} 

\begin{example} 
\rref{suppisclosed} \emph{salah} untuk modul yang tidak terbangkitkan
hingga. Sebagai contoh, pertimbangkan grup abelian $\bigoplus_p \mathbb{Z}/p$.
Dukungan objek ini sebagai modul-$\mathbb{Z}$ tepat merupakan himpunan semua
titik tertutup (yakni ideal maksimal) dari $\spec \mathbb{Z}$, sehingga
tidak tertutup. 
\end{example} 

\begin{corollary} 
Gelanggang $R$ mempunyai sejumlah hingga ideal prima minimal. 
\end{corollary} 
\begin{proof} 
Jelas bahwa $\supp R = \spec R$. Jadi, setiap ideal prima $R$
memuat prima terkait dari $R$ berdasarkan \rref{supportassociated}.
\end{proof} 

Dengan demikian, $\spec R$ merupakan gabungan hingga komponen tertutup tak tereduksi
$\overline{\mathfrak{q}}$ jika $R$ noetherian.
\begin{quote}\small\textbf{Jembatan edisi (asli).}
Suatu himpunan bagian tertutup tak kosong disebut tak tereduksi jika bukan gabungan dua himpunan bagian
tertutup sejati. Penutupan suatu titik dari $\spec R$ bersifat tak tereduksi.\end{quote}

Kita baru saja melihat bahwa $\supp M$ merupakan himpunan bagian tertutup dari $\spec R$ dan merupakan gabungan
sejumlah hingga himpunan bagian tak tereduksi. Lebih tepatnya, 
\[ \supp M = \bigcup_{\mathfrak{q} \in \ass(M)}
\overline{\left\{\mathfrak{q}\right\}}  \]
meskipun mungkin terdapat pengulangan dalam pernyataan ini. Suatu prima terkait mungkin termuat
dalam prima terkait lainnya.  

\begin{definition} 
Suatu ideal prima $\mathfrak{p} \in \ass(M)$ disebut prima terkait \textbf{terisolasi} dari
$M$ jika minimal (terhadap urutan pada $\ass(M)$); selain itu, ideal tersebut disebut
\textbf{tertanam}. 
\end{definition} 

Jadi, prima-prima tertanam tidak diperlukan untuk mendeskripsikan dukungan $M$. 

\begin{quote}\small\textbf{Jembatan edisi (asli).}
Misalkan $R=k[x,y]$ dan $M=R/(x^2,xy)$. Dekomposisi
\[
  (x^2,xy)=(x)\cap(x^2,y)
\]
menunjukkan bahwa $(x)$ merupakan prima terkait terisolasi dari $M$, sedangkan ideal
maksimal $(x,y)$ tertanam.\end{quote}

\begin{remark} 
Akibatnya, dalam suatu gelanggang noetherian, setiap ideal prima minimal terdiri atas
pembagi nol. Meskipun tidak akan digunakan kelak, pernyataan yang sama juga benar
dalam gelanggang non-noetherian. Berikut suatu argumen.

Misalkan $R$ suatu gelanggang dan $\mathfrak{p} \subset R$ suatu ideal prima minimal. Maka
$R_{\mathfrak{p}}$ mempunyai tepat satu ideal prima.
Sekarang kita gunakan:

\begin{lemma} 
Jika suatu gelanggang $R$ mempunyai tepat satu ideal prima $\mathfrak{p}$, maka setiap $x \in
\mathfrak{p}$ nilpoten.
\end{lemma} 
\begin{proof} 
Memang, cukup ditunjukkan bahwa $R_x = 0$ (\rref{nilpcriterion} dalam
\rref{spec}) jika $x \in
\mathfrak{p}$. Namun, $\spec R_x$
terdiri atas ideal-ideal prima $R$ yang tidak memuat $x$. Akan tetapi, tidak ada ideal prima semacam itu.
Jadi, $\spec R_x = \emptyset$, sehingga $R_x = 0$.
\end{proof} 

Akibatnya, setiap elemen dalam $\mathfrak{p}$ merupakan pembagi nol dalam
$R_{\mathfrak{p}}$.
Oleh karena itu, jika $x \in \mathfrak{p}$, terdapat $\frac{s}{t} \in
R_{\mathfrak{p}}$ sedemikian sehingga $xs/t = 0$, tetapi $\frac{s}{t} \neq 0$.
Secara khusus, terdapat $t' \notin \mathfrak{p}$ dengan 
\[ xst' = 0, \quad st' \neq 0,  \]
sehingga $x$ merupakan pembagi nol. 
\end{remark}



\subsection{Modul primer}

Modul primer adalah modul yang hanya mempunyai satu prima terkait. Hal ini setara
dengan mengatakan bahwa setiap homoteti bersifat injektif atau nilpoten. 
Seperti akan kita lihat pada bagian berikutnya, setiap modul mempunyai suatu “dekomposisi
primer”: bahkan, modul tersebut dapat disisipkan sebagai submodul dari suatu jumlah modul-modul
primer.

\begin{definition} 
Misalkan $\mathfrak{p} \subset R$ suatu ideal prima dan $M$ suatu modul-$R$ terbangkitkan hingga. Maka $M$ disebut
\textbf{$\mathfrak{p}$-primer} jika 
\[ \ass(M) = \left\{\mathfrak{p}\right\}.  \]

Suatu modul disebut \textbf{primer} jika $\mathfrak{p}$-primer untuk suatu
ideal prima $\mathfrak{p}$, yakni mempunyai tepat satu prima terkait. 
\end{definition} 

\begin{proposition} \label{whenisprimary}
Misalkan $M$ suatu modul-$R$ terbangkitkan hingga. Maka $M$ bersifat \textbf{$\mathfrak{p}$}-primer jika
dan hanya jika, untuk setiap $m \in M - \left\{0\right\}$, 
anihilator $\ann(m)$ mempunyai radikal $\mathfrak{p}$.
\end{proposition} 
\begin{proof} 
Pertama-tama kita memerlukan suatu pengamatan kecil.

\begin{lemma} 
Jika $M$ bersifat $\mathfrak{p}$-primer, maka setiap submodul tak nol $M' \subset M$ juga
bersifat $\mathfrak{p}$-primer.
\end{lemma} 
\begin{proof} 
Memang, kita mengetahui bahwa $\ass(M') \subset \ass(M)$ berdasarkan \rref{assexact}.
Karena $M' \neq 0$, kita juga mengetahui bahwa $M'$ mempunyai prima terkait
(\rref{assmnonempty}). Jadi, $ \ass(M') = \{\mathfrak{p}\}$, sehingga
$M'$ bersifat $\mathfrak{p}$-primer.
\end{proof} 

Sekarang kita kembali ke pembuktian hasil utama,
\rref{whenisprimary}.
Pertama, andaikan $M$ bersifat $\mathfrak{p}$-primer. Misalkan $x \in M$, $x \neq 0$. Tetapkan
$I = \ann(x)$; kita harus menunjukkan bahwa $\rad(I)  =\mathfrak{p}$. Berdasarkan definisi, terdapat penyisipan
\[ R/I \hookrightarrow M  \]
yang memetakan $1 \to x$. Akibatnya, $R/I$ bersifat $\mathfrak{p}$-primer berdasarkan lema di atas.
Kita ingin mengetahui bahwa $\mathfrak{p}  = \rad(I)$. 
Kita telah melihat bahwa dukungan $\supp R/I = \left\{\mathfrak{q}: \mathfrak{q}
\supset I\right\}$ merupakan gabungan penutupan prima terkait. Dalam
kasus ini, 
\[ \supp(R/I) = \left\{\mathfrak{q}: \mathfrak{q} \supset \mathfrak{p}\right\}
.\]
Namun, kita mengetahui bahwa $\rad(I) = \bigcap_{\mathfrak{q} \supset I} \mathfrak{q}$,
yang berdasarkan keterangan di atas tepat sama dengan $\mathfrak{p}$. Ini membuktikan bahwa $\rad(I) =
\mathfrak{p}$.
Kita telah menunjukkan bahwa jika $R/I$ primer, maka $I$ mempunyai radikal $\mathfrak{p}$.

Arah sebaliknya mudah. 
Andaikan syarat tersebut berlaku dan $\mathfrak{q} \in \ass(M)$, sehingga $\mathfrak{q} =
\ann(x)$ untuk $x \neq 0$. Namun, $\rad(\mathfrak{q}) = \mathfrak{p}$, sehingga 
\[ \mathfrak{q} = \mathfrak{p}  \] dan $\ass(M) = \left\{\mathfrak{p}\right\}$.
\end{proof} 

Kita mempunyai pencirian lain.

\begin{proposition} \label{whenisprimary2}
Misalkan $M \neq 0$ suatu modul-$R$ terbangkitkan hingga. Maka $M$ primer jika dan
hanya jika untuk setiap $a \in
R$, homoteti $ M \stackrel{a}{\to} M$ bersifat injektif atau nilpoten.
\end{proposition}
\begin{proof} 
Pertama, andaikan $M$ bersifat $\mathfrak{p}$-primer. Perkalian dengan setiap elemen
$\mathfrak{p}$ bersifat nilpoten karena anihilator setiap elemen tak nol mempunyai
radikal $\mathfrak{p}$ berdasarkan \rref{whenisprimary}. Namun, jika $a \notin \mathfrak{p}$, maka $\ann(x)$ untuk
$x \in M - \left\{0\right\}$ mempunyai radikal $\mathfrak{p}$ dan tidak mungkin memuat $a$. 

Sekarang kita buktikan arah sebaliknya. Andaikan setiap elemen $a$ bekerja secara injektif atau nilpoten pada $M$.
Misalkan $I \subset R$ kumpulan elemen $a \in R$ sedemikian sehingga $a^n M = 0$
untuk $n$ yang cukup besar. Maka $I$ merupakan ideal, sebab tertutup terhadap penjumlahan berdasarkan
rumus binomial: jika $a, b \in I$ serta $a^n, b^n$ bekerja sebagai nol, maka $(a+b)^{2n}$
juga bekerja sebagai nol.


Kita mengklaim bahwa $I$ sebenarnya prima. Jika $a,b \notin I$, maka $a,b$ bekerja melalui
perkalian secara injektif pada $M$. Jadi, $a: M \to M, b: M \to M$ injektif.
Namun, komposisi pemetaan injektif bersifat injektif, sehingga $ab$ bekerja secara injektif dan
$ab \notin I$. Jadi, $I$ prima.

Sekarang kita perlu memeriksa bahwa jika $x \in M$ tak nol, maka $\ann(x)$ mempunyai radikal
$I$. Memang, jika $a \in R$ menganihilasi $x$,
maka homoteti $M \stackrel{a}{\to} M$ tidak mungkin injektif, sehingga harus
nilpoten (yakni termasuk dalam $I$). Sebaliknya, jika $a \in I$, maka suatu pangkat $a$
bersifat nilpoten, sehingga suatu pangkat $a$ 
harus membunuh $x$. 
Akibatnya, $\ann(x) = I$. Sekarang, berdasarkan \rref{whenisprimary}, kita melihat
bahwa $M$ bersifat $I$-primer.
\end{proof} 

Sekarang kita mempunyai konsep modul primer. Gagasannya ialah bahwa seluruh torsi
dalam suatu cara terkonsentrasi pada suatu ideal prima.

\begin{example} 
Jika $R$ suatu gelanggang noetherian dan $\mathfrak{p} \in \spec R$, maka
$R/\mathfrak{p}$ bersifat $\mathfrak{p}$-primer. Secara lebih umum, jika $I \subset R$
suatu ideal, maka $R/I$ primer jika dan hanya jika $\rad(I) $ prima. Hal ini
mengikuti dari \rref{whenisprimary2}.
\end{example} 

\begin{exercise} 
Jika $0 \to M' \to M \to M'' \to 0$ suatu barisan eksak dengan $M', M, M''$
tak nol dan terbangkitkan hingga, maka $M$ bersifat $\mathfrak{p}$-primer jika dan hanya jika $M', M''$ demikian.
\end{exercise} 

\begin{exercise} 
Misalkan $M$ suatu modul-$R$ terbangkitkan hingga. Misalkan $\mathfrak{p} \in \spec R$. Tunjukkan bahwa jumlah dua
submodul $\mathfrak{p}$-primer juga $\mathfrak{p}$-primer. Simpulkan bahwa
terdapat suatu submodul $\mathfrak{p}$-primer dari $M$ yang memuat setiap
submodul $\mathfrak{p}$-primer. 
\end{exercise} 

\begin{exercise}[Bourbaki]
Misalkan $M$ suatu modul-$R$ terbangkitkan hingga. Misalkan $T \subset \ass(M)$ suatu
himpunan bagian dari prima terkait. Buktikan bahwa terdapat submodul $N \subset M$
sedemikian sehingga 
\[ \ass(N) = T, \quad \ass(M/N) = \ass(M) - T.  \]

\end{exercise} 

\section{Dekomposisi primer} Dalam suatu pengertian, inilah teorema struktur untuk modul
di atas gelanggang noetherian.
Di seluruh pembahasan, kita tetapkan suatu gelanggang noetherian $R$.

\subsection{Modul tak tereduksi dan koprimer}

\begin{definition} 
Misalkan $M$ suatu modul-$R$ terbangkitkan hingga. Suatu submodul $N \subset M$ disebut
\textbf{$\mathfrak{p}$-koprimer} jika $M/N$ bersifat $\mathfrak{p}$-primer.

Demikian pula, kita dapat mengatakan bahwa $N \subset M$ bersifat \textbf{koprimer} jika
$\mathfrak{p}$-koprimer untuk suatu $\mathfrak{p} \in \spec R$.
\end{definition} 

Sekarang kita akan menunjukkan bahwa setiap submodul $M$ dapat dinyatakan sebagai irisan
submodul-submodul koprimer. Untuk melakukannya, kita definisikan suatu submodul $M$ sebagai
\emph{tak tereduksi} jika tidak dapat ditulis sebagai irisan tak trivial dari
submodul-submodul $M$. Berdasarkan penalaran umum,
setiap submodul akan merupakan irisan submodul-submodul
tak tereduksi. Kemudian kita akan melihat bahwa setiap submodul tak tereduksi bersifat
koprimer. 

\begin{definition} 
Submodul $N \subsetneq M$ disebut \textbf{tak tereduksi} jika, setiap kali $N = N_1 \cap N_2$ untuk submodul $N_1,
N_2 \subset M$, salah satu dari $N_1, N_2$ sama dengan $N$. Dengan kata lain, submodul ini bukan
 irisan submodul-submodul yang lebih besar. 
\end{definition} 

\begin{proposition} \label{irrediscoprimary}
Suatu submodul tak tereduksi $N \subset M$ bersifat koprimer.
\end{proposition} 
\begin{proof} 
Misalkan $a \in R$. Kita ingin menunjukkan bahwa homoteti 
\[ M/N \stackrel{a}{\to} M/N  \]
bersifat injektif atau nilpoten.
Pertimbangkan submodul-submodul $M/N$ berikut:
\[ K(n) =  \left\{x \in M/N: a^n x = 0\right\} . \]
Jelas bahwa $K(0) \subset K(1) \subset \dots$; rantai ini menjadi stasioner karena
modul hasil bagi tersebut noetherian.
Secara khusus, $K(n) = K(2n)$ untuk $n$ yang cukup besar. 

Akibatnya, jika $x \in M/N$ habis dibagi $a^n$ ($n$ cukup besar) dan tak nol, maka $a^n x$
juga tak nol. Memang, tuliskan $x = a^n y \neq 0$; maka $y \notin K(n)$, sehingga $a^{n}x =
a^{2n}y \neq 0$, sebab jika tidak, kita memperoleh $y \in K(2n) = K(n)$. Dalam $M/N$, submodul-submodul
\[ a^n(M/N) \cap \ker(a^n)  \]
sama dengan nol untuk $n$ yang cukup besar. Namun, asumsi kita adalah bahwa $N$
tak tereduksi. Jadi, salah satu dari submodul-submodul $M/N$ ini nol. Artinya, berlaku
$a^n(M/N) = 0$ atau $\ker a^n = 0$. Kita memperoleh sifat injektif atau sifat nilpoten pada
$M/N$. Ini membuktikan hasil tersebut.
\end{proof} 
 
\subsection{Dekomposisi tak tereduksi dan dekomposisi primer}

Sekarang kita akan menunjukkan bahwa dalam modul terbangkitkan hingga di atas gelanggang noetherian,
kita dapat menulis $0$ sebagai irisan modul-modul koprimer. Dekomposisi ini,
yang disebut \emph{dekomposisi primer}, akan diturunkan dari penalaran yang sepenuhnya
umum.

\begin{definition} 
Suatu \textbf{dekomposisi tak tereduksi} dari modul $M$ adalah penyajian
$N_1 \cap N_2 \dots \cap N_k  = 0$, dengan $N_i \subset M$ merupakan submodul-submodul tak tereduksi
dalam dekomposisi tersebut.
\end{definition} 

\begin{proposition} 
Jika $M$ terbangkitkan hingga, maka $M$ mempunyai dekomposisi tak tereduksi. Terdapat sejumlah hingga submodul tak tereduksi
$N_1, \dots, N_k$ dengan
\[  N_1 \cap \dots \cap N_k = 0. \]
\end{proposition} 
Dengan kata lain,
\[  M \to \bigoplus M/N_i  \]
bersifat injektif.
Jadi, suatu modul terbangkitkan hingga di atas gelanggang noetherian dapat disisipkan ke dalam jumlah langsung
modul-modul primer, sebab berdasarkan \rref{irrediscoprimary}, semua $M/N_i$
bersifat primer. 

\begin{proof} Sekarang pembuktian sepenuhnya formal. 

Di antara submodul-submodul $M$, beberapa mungkin dapat dinyatakan sebagai irisan
sejumlah hingga submodul tak tereduksi, sedangkan yang lain mungkin tidak. Tujuan kita ialah menunjukkan bahwa
$0$ merupakan irisan semacam itu.
Misalkan $M' \subset M$ suatu submodul maksimal dari $M$ sedemikian sehingga $M'$ \emph{tidak dapat}
ditulis sebagai irisan semacam itu. Jika tidak ada
$M'$ seperti itu, pembuktian selesai, sebab $0$ dapat ditulis sebagai
irisan sejumlah hingga submodul tak tereduksi.

Sekarang $M'$ tidak tak tereduksi, sebab jika demikian, ia merupakan irisan satu submodul
tak tereduksi.
Akibatnya, $M'$ dapat ditulis sebagai $M'=M_1' \cap M_2'$ untuk dua submodul
$M$ yang lebih besar secara ketat. Namun, berdasarkan kemaksimalan, $M_1', M_2'$ mempunyai dekomposisi sebagai
irisan submodul-submodul tak tereduksi. Jadi, $M'$ juga mempunyai dekomposisi semacam itu, suatu kontradiksi. 
\end{proof} 

\begin{corollary} 
Untuk setiap $M$ terbangkitkan hingga, terdapat submodul-submodul koprimer $N_1, \dots,
N_k \subset M$ sedemikian sehingga $N_1 \cap \dots \cap N_k  = 0$.
\end{corollary} 
\begin{proof} 
Memang, setiap submodul tak tereduksi bersifat koprimer.
\end{proof} 


Untuk setiap $M$, kita mempunyai \textbf{dekomposisi tak tereduksi}
\[ 0 = \bigcap N_i  \]
dengan $N_i$ suatu himpunan hingga submodul tak tereduksi (dan dengan demikian koprimer). 
Dekomposisi ini sangat tidak tunggal dan tidak kanonik. Mari kita
sederhanakan menjadi sesuatu yang jauh lebih kanonik.

Klaim pertama adalah bahwa kita dapat mengelompokkan modul-modul yang koprimer terhadap
ideal prima yang sama. 
\begin{lemma} 
Misalkan $N_1, N_2 \subset M$ submodul-submodul $\mathfrak{p}$-koprimer. Maka $N_1 \cap
N_2$ juga $\mathfrak{p}$-koprimer.
\end{lemma} 
\begin{proof} 
Kita harus menunjukkan bahwa $M/N_1 \cap N_2$ bersifat $\mathfrak{p}$-primer. Memang, kita mempunyai penyisipan
\[ M/N_1 \cap N_2 \rightarrowtail  M/N_1 \oplus M/N_2  \]
yang menyiratkan bahwa $\ass(M/N_1 \cap N_2) \subset \ass(M/N_1) \cup \ass(M/N_2) =
\left\{\mathfrak{p}\right\}$. Jadi, pembuktian selesai. 
\end{proof} 

Secara khusus, jika kita tidak mensyaratkan sifat tak tereduksi dan hanya mensyaratkan sifat primer dalam
dekomposisi
\[ 0 = \bigcap N_i,  \]
kita dapat mengandaikan bahwa setiap $N_i$ bersifat $\mathfrak{p}_i$-koprimer untuk suatu
 ideal prima
$\mathfrak{p}_i$, dengan semua $\mathfrak{p}_i$ \emph{berbeda}.

\begin{definition} 
Dekomposisi nol semacam itu, dengan modul-modul $N_i$ yang berbeda bersifat
$\mathfrak{p}_i$-koprimer untuk $\mathfrak{p}_i$ yang berbeda, disebut suatu \textbf{dekomposisi primer}.
\end{definition} 



\subsection{Pertanyaan ketunggalan}

Secara umum, dekomposisi primer \emph{tidak} tunggal. Meskipun demikian, kita akan
melihat bahwa terdapat ketunggalan dalam kadar terbatas. Sebagai contoh, ideal-ideal prima
yang muncul ditentukan secara tunggal.

Misalkan $M$ suatu modul terbangkitkan hingga di atas gelanggang noetherian $R$, dan andaikan
$N_1 \cap \dots \cap N_k = 0$ suatu dekomposisi primer.
Kita andaikan bahwa dekomposisi tersebut
\emph{minimal}: artinya, jika salah satu $N_i$ dihapus, irisannya tidak lagi
nol.
Hal ini menyiratkan bahwa
\[ N_i \not\supset \bigcap_{j \neq i} N_j  \]
sebab jika tidak, salah satu $N_i$ dapat dihilangkan. Dekomposisi tersebut lalu disebut \textbf{dekomposisi primer tereduksi}.

Sekali lagi, hal ini menyatakan bahwa $M \rightarrowtail  \bigoplus M/N_i$. Kita telah
menunjukkan bahwa $M$ dapat disisipkan ke dalam jumlah komponen, yang masing-masing
bersifat $\mathfrak{p}$-primer untuk suatu ideal prima, dan semua ideal prima tersebut berbeda.

Penyisipan ini \textbf{tidak} tunggal. Akan tetapi, 

\begin{proposition} 
Ideal-ideal prima $\mathfrak{p}_i$ yang muncul dalam suatu dekomposisi primer tereduksi dari nol
ditentukan secara tunggal. Ideal-ideal tersebut adalah prima terkait dari $M$.
\end{proposition} 
\begin{proof} 
Semua prima terkait dari $M$ harus muncul, karena kita mempunyai penyisipan
\[ M \rightarrowtail  \bigoplus M/N_i  \]
sehingga prima terkait dari $M$ termasuk di antara prima terkait dari $M/N_i$ (yakni semua
$\mathfrak{p}_i$).

Arah yang sulit adalah menunjukkan bahwa setiap $\mathfrak{p}_i$ merupakan prima terkait.
Artinya, jika $M/N_i$ bersifat $\mathfrak{p}_i$-primer, maka $\mathfrak{p}_i \in \ass(M)$;
kita tidak perlu menggunakan modul primer kecuali bagi ideal prima dalam prima terkait. 
Di sini kita perlu menggunakan fakta bahwa dekomposisi kita tidak berulang. Tanpa
mengurangi keumuman, cukup ditunjukkan bahwa, misalnya, $\mathfrak{p}_1$
termasuk dalam $\ass(M)$. Kita akan menggunakan fakta bahwa
\[ N_1 \not\supset N_2 \cap \dots .  \]
Jadi, $N_2 \cap N_3 \cap \dots$ tidak sama dengan nol, sebab jika sama,
kita akan memperoleh suatu penyertaan. Kita mempunyai pemetaan
\[ N_2 \cap \dots \cap N_k \to M/N_1;  \]
yang injektif, sebab kernelnya adalah $N_1 \cap N_2 \cap \dots \cap N_k = 0$ karena
ini merupakan suatu dekomposisi.
Namun, $M/N_1$ bersifat $\mathfrak{p}_1$-primer, sehingga $N_2 \cap \dots \cap N_k$ juga
bersifat $\mathfrak{p}_1$-primer. Secara khusus, $\mathfrak{p}_1$ merupakan prima terkait
dari $N_2 \cap \dots \cap N_k$, dan karena itu juga dari $M$.
\end{proof} 

Ideal-ideal primanya ditentukan secara tunggal; faktornya tidak. Namun, dalam beberapa kasus, faktornya tunggal.

\begin{proposition} 
Misalkan $\mathfrak{p}_i$ suatu prima terkait minimal dari $M$, yakni tidak memuat
prima terkait yang lebih kecil. Maka submodul $N_i$ yang bersesuaian dengan
$\mathfrak{p}_i$ dalam dekomposisi primer tereduksi ditentukan secara tunggal:
submodul tersebut adalah kernel dari 
\[ M \to M_{\mathfrak{p}_i}.  \]
\end{proposition} 

\begin{proof} 
Kita mempunyai $\bigcap N_j = \left\{0\right\} \subset M$. Ketika dilokalkan pada
$\mathfrak{p}_i$, kita memperoleh
\[ (\bigcap N_j)_{\mathfrak{p}_i} = \bigcap (N_j)_{\mathfrak{p}_i} =0 \]
karena pelokalan merupakan funktor eksak. Jika $j \neq i$, maka $M/N_j$ bersifat
$\mathfrak{p}_j$-primer dan hanya mempunyai $\mathfrak{p}_j$ sebagai prima terkait.
Akibatnya, $(M/N_j)_{\mathfrak{p}_i}$ tidak mempunyai prima terkait, karena satu-satunya
prima terkait yang mungkin adalah $\mathfrak{p}_j$, dan ideal tersebut tidak termuat dalam
$\mathfrak{p}_j$.
Secara khusus, $(N_j)_{\mathfrak{p}_i} = M_{\mathfrak{p}_i}$.

Jadi, ketika kita melokalkan dekomposisi primer pada $\mathfrak{p}_i$, kita memperoleh
dekomposisi primer trivial: sebagian besar faktornya adalah seluruh
$M_{\mathfrak{p}_i}$. Akibatnya, $(N_i)_{\mathfrak{p}_i}=0$. Ketika kita menggambar
diagram komutatif
\[ 
\xymatrix{
N_i \ar[r] \ar[d]  &  (N_i)_{\mathfrak{p}_i} = 0 \ar[d]  \\
M \ar[r] &  M_{\mathfrak{p}_i}.
}
\]
kita mendapati bahwa $N_i$ dipetakan ke nol dalam pelokalan.

Sekarang, jika $x \in \ker(M \to M_{\mathfrak{p}_i}$, maka $sx = 0$ untuk suatu $s \notin
\mathfrak{p}_i$. Ketika kita mengambil pemetaan $M \to M/N_i$, $sx$ dipetakan ke nol; tetapi $s$
bekerja secara injektif pada $M/N_i$, sehingga $x$ dipetakan ke nol di $M/N_i$, yakni termasuk dalam
$N_i$.
\end{proof} 

Pembahasan ini bersifat abstrak, jadi pertimbangkan contoh berikut:
\begin{example} Misalkan $ R = \mathbb{Z}$.
Misalkan $M = \mathbb{Z} \oplus \mathbb{Z}/p$. Maka nol dapat ditulis sebagai 
\[ \mathbb{Z} \cap \mathbb{Z}/p  \]
sebagai submodul-submodul $M$. Namun, $\mathbb{Z}$ bersifat $\mathfrak{p}$-koprimer, sedangkan
$\mathbb{Z}/p$ bersifat $(0)$-koprimer. 

Dekomposisi ini tidak tunggal. Kita dapat mempertimbangkan 
\[ \{(n,n), n \in \mathbb{Z}\} \subset M.  \]
Namun, bagian nol-koprimer haruslah torsi-$p$. Hal ini karena
$(0)$ merupakan ideal minimal. 

Secara umum, dekomposisi selalu tunggal jika
tidak terdapat relasi penyertaan di antara ideal-ideal prima. Untuk modul-$\mathbb{Z}$, ini
berarti bahwa dekomposisi primer tunggal untuk modul torsi. 
Setiap grup torsi merupakan jumlah langsung dari torsi berpangkat-$p$ untuk semua bilangan prima $p$. 
\end{example} 

\begin{exercise} 
Andaikan $R$ suatu gelanggang noetherian dan $R_{\mathfrak{p}}$ suatu domain untuk setiap ideal prima
$\mathfrak{p} \subset R$. Maka $R$ merupakan hasil kali langsung hingga $\prod R_i$, dengan
setiap $R_i$ suatu domain.

Untuk melihatnya, pertimbangkan ideal-ideal prima minimal $\mathfrak{p}_i \in \spec R$. Terdapat
sejumlah hingga ideal semacam itu; buktikan bahwa karena setiap pelokalan merupakan domain,
$\spec R$ terpisah menjadi komponen-komponen $V(\mathfrak{p}_i)$.
Akibatnya, terdapat dekomposisi $R = \prod R_{i}$ dengan $\spec R_i$
mempunyai $\mathfrak{p}_i$ sebagai satu-satunya ideal prima minimal. 
Setiap $R_i$ memenuhi syarat yang sama seperti $R$, sehingga kita dapat mereduksi ke kasus
ketika $R$ mempunyai satu-satunya ideal prima minimal. Namun, dalam kasus ini $R$
tereduksi, sehingga satu-satunya ideal prima minimal haruslah nol.
\end{exercise} 



